\section{Challenges of Context Retrieval for Coding Agents}
\label{sec:ch1_bg}

% \myx{Maybe this sub-section should be the first one? Also, should not be background, but something like "Challenges of context retrieval for coding agents."}

% \frank{this subsection does go through some design and implementation details on RepoGraph which I think should come later than the previous subsection, does it make sense if I simply move details on RepoGraph to the end of previous section then switch their order? The relevant parts are marked magenta below}

Code repositories present unique challenges for LLM-based agents due to their scale and interconnectivity. Real-world codebases often contain hundreds of thousands of lines of code spread across many files. This makes it infeasible to simply dump the entire repository into the LLM’s context window. Instead, an agent must \emph{localize} the relevant portions of code for a given task. Traditional program analysis techniques, such as dependency graphs and call graphs, have long been used to understand large codebases, and recent LLM agents are beginning to incorporate such ideas. Early coding assistants relied on heuristic searches (e.g.\ grepping for error messages or function names) and would open files one by one in a Localize-then-Edit paradigm. More modern agents improve on this by indexing repository knowledge: for example, SWE-Search \cite{antoniades2025swesearchenhancingsoftwareagents} integrate a Monte Carlo Tree Search to explore the complex solution space, and CodePlan \cite{bairi2023codeplanrepositorylevelcodingusing} tracks incremental code changes with dependency analysis.

Despite these advancements, providing just the right context remains difficult. A key issue is that repository-level issue solving often requires following chains of references across multiple files. If any link in the chain is missing from the prompt, the model may fail to resolve the problem. Conversely, adding too much code can dilute the context. CoSIL \cite{jiang2025cosilsoftwareissuelocalization} \myx{Use this to support the sentence "Conversely, naively supplying too much code ..." in Sec. 2.1.} highlights this tension: it points out that existing methods struggle to give concise yet effective context coverage . CoSIL’s solution is to dynamically construct a module-level call graph during the agent’s search, expanding the graph by following function calls and pruning irrelevant branches. RepoGraph takes a complementary approach by pre-building a complete repository graph before the agent starts working. During a setup phase, RepoGraph parses all code files (using \texttt{tree-sitter}) to identify definitions (classes, functions, variables) and their references, filters out non-code files, and stores the resulting graph for fast querying. When the agent needs context for a particular function or error, it can query this graph to retrieve all related definitions and usage chains. In essence, RepoGraph provides an oracle that, given a symbol or filename, returns a subgraph containing all immediately relevant symbols.
